% Source: https://artofproblemsolving.com/community/c3h2144867_inverse_functions
% Stable identity: aops-topic-2144867-q01
Since $f$ is the inverse of $g$, we have $f(4)=x$ exactly when $g(x)=4$. (This is the method suggested in the post, and it is valid: $g$ is one-to-one, so there is only one such $x$.) Solve
\[\dfrac{2x-3}{x+5}=4\iff 2x-3=4x+20\iff x=-\dfrac{23}{2}.\]
Check: $g\left(-\dfrac{23}{2}\right)=\dfrac{-23-3}{-\dfrac{23}{2}+5}=\dfrac{-26}{-\dfrac{13}{2}}=4$. So $f(4)=\boxed{-\dfrac{23}{2}}$.

\emph{Alternatively}, find $f$ itself: from $y=\dfrac{2x-3}{x+5}$ we get $xy+5y=2x-3$, so $x=\dfrac{5y+3}{2-y}$. Hence $f(x)=\dfrac{5x+3}{2-x}$ (for $x\ne2$), and $f(4)=\dfrac{23}{-2}=-\dfrac{23}{2}$.
